\subsection{Dependencias}

\begin{itemize}
  \item \texttt{make}.
  \item Toolchain bare-metal: \texttt{riscv64-unknown-elf-\{gcc, g++, as, ld, objcopy\}}.
  \item \texttt{qemu-system-riscv64}.
\end{itemize}

\subsection{Compilación}

\begin{codesnippet}
\begin{verbatim}
 make -C riscv clean all
\end{verbatim}
\end{codesnippet}

Genera \texttt{riscv/build/kernel.elf} (y \texttt{.bin}), que contiene embebidos los binarios de usuario
(\texttt{build/user\_miner.bin}, \texttt{build/user\_idle.bin}).

\subsection{Ejecución}

\begin{codesnippet}
\begin{verbatim}
 make -C riscv run
\end{verbatim}
\end{codesnippet}

que se expande a:

\begin{codesnippet}
\begin{verbatim}
 qemu-system-riscv64 -machine virt -nographic -monitor none \
     -serial stdio -bios none -kernel riscv/build/kernel.elf
\end{verbatim}
\end{codesnippet}

\subsection{Controles}

\begin{itemize}
  \item Jugador ROJO: \texttt{W} (arriba), \texttt{S} (abajo), \texttt{Q} o \texttt{E} (spawnear minero).
  \item Jugador AZUL: \texttt{I} (arriba), \texttt{K} (abajo), \texttt{O} o \texttt{P} (spawnear minero).
\end{itemize}

En la esquina superior derecha se muestra, en hex, el último byte recibido por el UART (útil para validar input).
